
\documentclass[11pt,a4paper]{article}
\usepackage[margin=0.78in,headheight=14pt,footskip=24pt]{geometry}
\usepackage[T1]{fontenc}
\usepackage{lmodern}
\usepackage{xcolor,listings,graphicx,float,fancyhdr,booktabs,caption,titlesec,microtype,array}
\definecolor{codebg}{RGB}{247,247,247}
\definecolor{codegreen}{RGB}{0,110,55}
\definecolor{codeblue}{RGB}{30,70,160}
\definecolor{codepurple}{RGB}{135,40,140}
\lstdefinestyle{terminalstyle}{
 basicstyle=\ttfamily\scriptsize, backgroundcolor=\color{codebg},
 breaklines=true, frame=single, rulecolor=\color{black!18},
 columns=fullflexible, keepspaces=true, showstringspaces=false
}
\lstdefinestyle{pythonstyle}{
 language=Python, basicstyle=\ttfamily\scriptsize,
 backgroundcolor=\color{codebg}, keywordstyle=\color{codeblue}\bfseries,
 commentstyle=\color{codegreen}\itshape, stringstyle=\color{codepurple},
 breaklines=true, frame=single, rulecolor=\color{black!18},
 numbers=left, numberstyle=\tiny\color{gray}, numbersep=6pt,
 columns=fullflexible, showstringspaces=false
}
\pagestyle{fancy}
\fancyhf{}
\lhead{\small Big Data Analysis Lab}
\rfoot{\small Experiment 10}
\cfoot{\thepage}
\renewcommand{\headrulewidth}{0.4pt}
\setlength{\parindent}{0pt}
\setlength{\parskip}{4pt}
\titleformat{\section}{\large\bfseries}{\thesection}{0.7em}{}
\titleformat{\subsection}{\normalsize\bfseries}{\thesubsection}{0.7em}{}
\captionsetup{font=small,labelfont=bf}
\setlength{\intextsep}{5pt}
\setlength{\textfloatsep}{6pt}
\setlength{\abovecaptionskip}{3pt}
\setlength{\belowcaptionskip}{0pt}
\begin{document}
\noindent
\textbf{Name:} Nikhil Sunchu\\
\textbf{Class:} D20A-31

\vspace{0.5em}
\begin{center}
{\Large\bfseries Experiment 10: E-Commerce Intelligence Dashboard}\\
\textbf{Analytics using Hive, Impala, PostgreSQL, and Streamlit}
\end{center}

\section{Implementation Steps and Results}

\subsection{1. Store the Dataset in HDFS}
The October 2019 e-commerce events CSV file (\texttt{2019-Oct.csv}, approximately 3.13 GB) was copied into the Cloudera QuickStart container and uploaded to HDFS. The file was stored in \texttt{/user/cloudera/exp10}.

\begin{lstlisting}[style=terminalstyle]
hdfs dfs -mkdir -p /user/cloudera/exp10
hdfs dfs -put /tmp/2019-Oct.csv /user/cloudera/exp10/
hdfs dfs -ls -h /user/cloudera/exp10
\end{lstlisting}

\subsection{2. Create the Hive External Table}
An external Hive table was created over the CSV file. The schema contains nine columns. The table uses comma-separated fields and skips the header row.

\begin{lstlisting}[style=terminalstyle]
CREATE EXTERNAL TABLE IF NOT EXISTS ecommerce_events (
    event_time STRING,
    event_type STRING,
    product_id BIGINT,
    category_id BIGINT,
    category_code STRING,
    brand STRING,
    price DOUBLE,
    user_id BIGINT,
    user_session STRING
)
ROW FORMAT DELIMITED
FIELDS TERMINATED BY ','
STORED AS TEXTFILE
LOCATION '/user/cloudera/exp10'
TBLPROPERTIES ("skip.header.line.count"="1");
\end{lstlisting}

The table was moved into \texttt{exp10\_db} and Impala metadata was refreshed.

\begin{lstlisting}[style=terminalstyle]
ALTER TABLE default.ecommerce_events
RENAME TO exp10_db.ecommerce_events;

INVALIDATE METADATA exp10_db.ecommerce_events;
SHOW TABLES IN exp10_db;
\end{lstlisting}

\begin{figure}[H]
\centering
\includegraphics[width=0.72\linewidth]{image(20261104-051210).png}
\caption{The \texttt{ecommerce\_events} table is visible in \texttt{exp10\_db} in Impala.}
\end{figure}

\subsection{3. Create a Clean Impala View}
Impala initially counted the CSV header as one additional record. A view was created to retain only the three valid event types: \texttt{view}, \texttt{cart}, and \texttt{purchase}.

\begin{lstlisting}[style=terminalstyle]
CREATE VIEW exp10_db.ecommerce_events_clean AS
SELECT *
FROM exp10_db.ecommerce_events
WHERE event_type IN ('view', 'cart', 'purchase');
\end{lstlisting}

The cleaned record count was checked using:

\begin{lstlisting}[style=terminalstyle]
SELECT COUNT(*) AS total_records
FROM exp10_db.ecommerce_events_clean;
\end{lstlisting}

The view returned 23,445,441 records, matching the Hive count.

\begin{figure}[H]
\centering
\includegraphics[width=0.82\linewidth]{image(20261104-051931).png}
\caption{Creation of the clean view and event distribution without the CSV header row.}
\end{figure}
\begin{figure}[H]
\centering
\includegraphics[width=0.68\linewidth]{image(20261104-052211).png}
\caption{Cleaned Impala record count: 23,445,441.}
\end{figure}

\subsection{4. Perform Analytics Using Impala}
The cleaned view was used for the following analytical queries.

\textbf{Event distribution:}
\begin{lstlisting}[style=terminalstyle]
SELECT event_type, COUNT(*) AS total_events
FROM exp10_db.ecommerce_events_clean
GROUP BY event_type
ORDER BY total_events DESC;
\end{lstlisting}
The results were 22,541,429 views, 487,353 cart events, and 416,659 purchases.

\textbf{Top 10 categories by views:}
\begin{lstlisting}[style=terminalstyle]
SELECT category_code, COUNT(*) AS total_views
FROM exp10_db.ecommerce_events_clean
WHERE event_type = 'view'
  AND category_code IS NOT NULL
  AND TRIM(category_code) != ''
GROUP BY category_code
ORDER BY total_views DESC
LIMIT 10;
\end{lstlisting}
\begin{figure}[H]
\centering
\includegraphics[width=0.72\linewidth]{image(20261104-052317).png}
\caption{Top 10 product categories by view events.}
\end{figure}

\textbf{Daily purchase trends:}
\begin{lstlisting}[style=terminalstyle]
SELECT SUBSTR(event_time, 1, 10) AS purchase_date,
       COUNT(*) AS purchase_events,
       ROUND(SUM(price), 2) AS daily_purchase_value
FROM exp10_db.ecommerce_events_clean
WHERE event_type = 'purchase'
GROUP BY SUBSTR(event_time, 1, 10)
ORDER BY purchase_date;
\end{lstlisting}
The query returned 18 dates from October 1 to October 18, 2019. October 18 contained 1,084 purchase events and is treated as a partial day.
\begin{figure}[H]
\centering
\includegraphics[width=0.78\linewidth]{image(20261104-052532).png}
\caption{Daily purchase-event counts and recorded purchase values.}
\end{figure}

\subsection{5. Create PostgreSQL Tables}
A PostgreSQL database named \texttt{ecommerce\_dashboard} was used to store the aggregated analytics. Three tables were created.

\begin{lstlisting}[style=terminalstyle]
CREATE TABLE event_summary (
    event_type VARCHAR(50) PRIMARY KEY,
    total_events BIGINT NOT NULL
);

CREATE TABLE daily_purchase_trends (
    purchase_date DATE PRIMARY KEY,
    purchase_events BIGINT NOT NULL,
    daily_purchase_value NUMERIC(15, 2)
);

CREATE TABLE category_performance (
    category_code VARCHAR(255) PRIMARY KEY,
    total_views BIGINT DEFAULT 0,
    purchase_events BIGINT DEFAULT 0,
    total_purchase_value NUMERIC(15, 2) DEFAULT 0
);
\end{lstlisting}
\begin{figure}[H]
\centering
\includegraphics[width=0.72\linewidth]{image(20261104-053116).png}
\caption{The three analytics tables created in PostgreSQL.}
\end{figure}

\subsection{6. Export Analytics from Impala to PostgreSQL}
The Python script \texttt{scripts/export\_analytics.py} connects to Impala using ImPyla on port 21050, extracts aggregated data into Pandas DataFrames, and inserts the results into PostgreSQL using SQLAlchemy.

\begin{lstlisting}[style=pythonstyle]
from impala.dbapi import connect
from sqlalchemy import create_engine, text
from dotenv import load_dotenv
import pandas as pd
import os

load_dotenv()

impala_conn = connect(
    host="172.22.0.2", port=21050,
    database="exp10_db", auth_mechanism="NOSASL"
)
cursor = impala_conn.cursor()

DATABASE_URL = (
    f"postgresql+psycopg2://{os.getenv('PGUSER')}:"
    f"{os.getenv('PGPASSWORD')}@{os.getenv('PGHOST')}:"
    f"{os.getenv('PGPORT')}/{os.getenv('PGDATABASE')}"
)
pg_engine = create_engine(DATABASE_URL)

def fetch_dataframe(query):
    cursor.execute(query)
    rows = cursor.fetchall()
    columns = [col[0] for col in cursor.description]
    return pd.DataFrame(rows, columns=columns)

event_df = fetch_dataframe("""
    SELECT event_type, COUNT(*) AS total_events
    FROM ecommerce_events_clean
    GROUP BY event_type
""")

daily_df = fetch_dataframe("""
    SELECT SUBSTR(event_time, 1, 10) AS purchase_date,
           COUNT(*) AS purchase_events,
           ROUND(SUM(price), 2) AS daily_purchase_value
    FROM ecommerce_events_clean
    WHERE event_type = 'purchase'
    GROUP BY SUBSTR(event_time, 1, 10)
    ORDER BY purchase_date
""")
daily_df["purchase_date"] = pd.to_datetime(
    daily_df["purchase_date"]
).dt.date

category_df = fetch_dataframe("""
    SELECT category_code,
           SUM(CASE WHEN event_type = 'view' THEN 1 ELSE 0 END)
               AS total_views,
           SUM(CASE WHEN event_type = 'purchase' THEN 1 ELSE 0 END)
               AS purchase_events,
           ROUND(SUM(CASE WHEN event_type = 'purchase'
                          THEN price ELSE 0 END), 2)
               AS total_purchase_value
    FROM ecommerce_events_clean
    WHERE category_code IS NOT NULL
      AND TRIM(category_code) != ''
    GROUP BY category_code
""")

with pg_engine.begin() as connection:
    connection.execute(text("""
        TRUNCATE TABLE event_summary,
        daily_purchase_trends, category_performance
    """))
    event_df.to_sql("event_summary", connection,
                    if_exists="append", index=False)
    daily_df.to_sql("daily_purchase_trends", connection,
                    if_exists="append", index=False)
    category_df.to_sql("category_performance", connection,
                       if_exists="append", index=False)

cursor.close()
impala_conn.close()
pg_engine.dispose()
\end{lstlisting}

The script exported 3 event-summary rows, 18 daily-trend rows, and 125 category rows.
\begin{figure}[H]
\centering
\includegraphics[width=0.78\linewidth]{image(20261104-064115).png}
\caption{PostgreSQL event-summary records.}
\end{figure}
\begin{figure}[H]
\centering
\includegraphics[width=0.78\linewidth]{image(20261104-064133).png}
\caption{PostgreSQL daily purchase-trend records.}
\end{figure}
\begin{figure}[H]
\centering
\includegraphics[width=0.94\linewidth]{image(20261104-064153).png}
\caption{Category-performance records stored in PostgreSQL.}
\end{figure}

\subsection{7. Build and Run the Streamlit Dashboard}
The Streamlit application in \texttt{dashboard/app.py} reads the three PostgreSQL tables and displays key performance indicators and Plotly charts. The dashboard was launched using:

\begin{lstlisting}[style=terminalstyle]
uv run streamlit run dashboard/app.py
\end{lstlisting}

The dashboard displays total events, product views, add-to-cart events, purchases, event distribution, daily purchase trends, daily recorded purchase value, and the top categories by views and purchases.

\begin{figure}[H]
\centering
\includegraphics[width=0.98\linewidth]{image(20261104-070129).png}
\caption{Dashboard KPIs, event distribution, daily purchase events, and data-coverage notice.}
\end{figure}
\begin{figure}[H]
\centering
\includegraphics[width=0.98\linewidth]{image(20261104-070150).png}
\caption{Daily recorded purchase value and top 10 categories by views.}
\end{figure}
\begin{figure}[H]
\centering
\includegraphics[width=0.98\linewidth]{image(20261104-070216).png}
\caption{Top 10 categories by purchase events and category analytics table.}
\end{figure}

\subsection{8. Verify the Final PostgreSQL Data}
The row counts in all three PostgreSQL tables were checked after the export.

\begin{lstlisting}[style=terminalstyle]
SELECT
    (SELECT COUNT(*) FROM event_summary) AS event_rows,
    (SELECT COUNT(*) FROM daily_purchase_trends) AS daily_rows,
    (SELECT COUNT(*) FROM category_performance) AS category_rows;
\end{lstlisting}

The final output was: \texttt{event\_rows = 3}, \texttt{daily\_rows = 18}, and \texttt{category\_rows = 125}. All three table counts matched the export summary.

\section{Observation}
\begin{center}
\begin{tabular}{>{\raggedright\arraybackslash}p{0.58\linewidth}p{0.32\linewidth}}
\toprule
\textbf{Operation} & \textbf{Observation} \\
\midrule
HDFS dataset & 2019-Oct.csv uploaded (approximately 3.13 GB) \\
Hive external table & Created in \texttt{exp10\_db} \\
Cleaned Impala view & 23,445,441 records \\
Event distribution & 3 event types returned \\
Daily purchase trends & 18 dates returned \\
Category analytics & 125 categories exported \\
PostgreSQL export & 3 event, 18 daily, 125 category rows \\
Streamlit dashboard & KPI cards and interactive charts displayed \\
Final validation & All PostgreSQL row counts matched \\
\bottomrule
\end{tabular}
\end{center}

\section{Conclusion}
The experiment successfully implemented an e-commerce analytics workflow using HDFS, Hive, Impala, Python, PostgreSQL, and Streamlit. The available dataset contained 23,445,441 valid event records and covered October 1--18, 2019. Aggregated event, daily purchase, and category analytics were transferred into PostgreSQL and visualised through the Streamlit dashboard. Final table counts were verified as 3, 18, and 125 rows respectively.

\end{document}
